% =====================================================================
%  ch3a_objects.tex — §3.1 对象与态射
% =====================================================================
\section{A 模型：三价构形中的拉格朗日 Floer 理论}\label{s:Amodel}

\subsection{$\F(M)$ 中的对象与态射}\label{ss:objects}

设 $G$ 为图，顶点集 $V$ 与边集 $E$ 均有限，且每个顶点 $v\in V$ 的度数为 $3$。如注~\ref{rmk:opencase} 所述，我们允许「外部边」：它只与一个顶点相连。记 $E^0$ 为外部边之集，$E^i$ 为内部边之集。对每条内部边 $e\in E^i$，固定面积参数 $A_e>0$ 与元素 $q_e\in K$，满足 $\val(q_e)=A_e$（我们主要关心 $q_e=T^{A_e}$ 的情形）；外部边没有面积参数，但我们考虑\emph{缠绕}或\emph{带停点}的拉格朗日 Floer 理论。

对每条内部边 $e\in E^i$，取 $M_e=S^2=\CP^1$，配备总面积为 $A_e$ 的辛形式 $\omega$（例如标准辛形式的一个倍数），并可配备体形变类 $\mathfrak{b}\in H^2(M_e,\cO_K)$，使得
\[
T^{A_e}\exp\Bigl(\int_{M_e}\mathfrak{b}\Bigr)=q_e .
\]
再在 $M_e$ 上固定两个标记点，视作 $\CP^1$ 中的 $0$ 与 $\infty$，并把它们分配给边 $e$ 所连接的两个顶点 $v,v'\in V$：
\[
\{p_{e/v},p_{e/v'}\}=\{0,\infty\}\subset M_e .
\]
对每条外部边 $e\in E^0$，令 $M_e=\C$，配以标准辛形式（无限面积）与单个标记点 $p_{e/v}=0\in M_e$。

设 $M$ 为把各曲面 $M_e$（$e\in E$）在对应于图 $G$ 同一顶点的三组标记点处互相粘合而得到的空间：
\[
M=\Bigl(\bigsqcup_{e\in E}M_e\Bigr)\Big/\bigl(p_{e/v}\sim p_{e'/v}\sim p_{e''/v}\ \ \forall v\in V\bigr).
\]
记所得的节点为 $p_v$。这一粘合纯粹是形式上的，真正的辛几何发生在各个分量 $M_e$ 上。另一方面，与每个顶点 $v\in V$ 相联系的一项关键数据，是抽象曲线 $X^0_v=\PP^1$ 上在相应标记点 $x_{e/v}\in X^0_v$ 处为零的局部坐标 $t_{e/v}$（比较定义~\ref{def:combidata}）。

在每个分量 $M_e\subset M$ 上，我们固定 $0$ 与 $\infty$ 附近的一个渐近方向，例如正实轴方向；所考虑的一切拉格朗日都必须沿这一指定方向趋近 $M$ 的节点与无穷远端。

\begin{defi}
$\F(M)$ 的\emph{对象}是配对 $(L,\cE)$，其中 $L\subset M$ 是 proper 嵌入的（三价）图，其顶点位于 $M$ 的节点处、边位于 $M$ 的光滑部分，并满足：
\begin{itemize}
\item $L_e=L\cap M_e$ 的弧分量沿指定渐近方向趋近 $0$ 与 $\infty$；
\item $L_e$ 的闭曲线分量在标记点补集中同伦非平凡；
\item 节点 $p_v\in M$ 属于 $L$，当且仅当它是交于 $p_v$ 的 $M$ 三个分量中各自一条弧的端点；
\end{itemize}
而 $\cE$ 是酉局部系统，即 $L$ 上自由有限秩 $\cO_K$-模的局部系统。
\end{defi}

由于 $M$ 的每个分量或是 $(\CP^1,\{0,\infty\})$ 或是 $(\C,\{0\})$，上述定义只允许两类不可分解对象。

\begin{enumerate}
\item \emph{点型对象}：$L$ 是某分量 $M_e$ 光滑部分中的简单闭曲线，分离 $0$ 与 $\infty$。当 $\cE$ 秩为 $1$ 时，对象 $(L,\cE)$ 在镜像对称下对应于 $X_K$ 上某点 $p$ 的摩天层，其中坐标 $t_{e/v}$ 在 $p$ 处的赋值等于 $L$ 绕标记点 $p_{e/v}$ 所围的辛面积。
\item \emph{向量丛（v.b.）型对象}：$L$ 是边集与顶点集均与 $G$ 相同的三价图，在每个分量 $M_e$ 中由一条连接 $0$ 与 $\infty$ 的弧 $L_e$ 组成，并经过所有节点。当 $\cE$ 秩为 $1$ 时，对象 $(L,\cE)$ 对应于 $X_K$ 上的一个线丛，详见 §\ref{ss:assignbundles}。
\end{enumerate}

我们还要指定一族光滑 Hamilton 扰动，用来定义 $\F(M)$ 各对象之间的 Floer 复形。

\begin{defi}
\emph{正 Hamilton 函数}是光滑函数 $h:M\to\R$，它在每个紧分量 $M_e\simeq\CP^1$（$e\in E^i$）上于两个标记点 $0,\infty$ 处取局部极小，$h(0)=h(\infty)=0$；在每个非紧分量 $M_e\simeq\C$（$e\in E^0$）上于原点取极小 $h(0)=0$，并且在无穷远处按坐标 $r=|z|^2$ 呈线性增长（当该非紧端带停点时）或二次增长（当不带停点时）。
\end{defi}

这样的 Hamilton 函数的流，把标记点附近的渐近方向沿正方向旋转，并以（部分）缠绕 Floer 理论的惯常方式推挤或缠绕无穷远端。

对每对 $(L,L')$，选取正 Hamilton 函数 $h$ 与足够小的 $\varepsilon>0$，使得 $L^+=\phi^1_{\varepsilon h}(L)$ 与 $L'$ 横截相交；定义 Floer 复形的生成元为由 $\varepsilon h$ 生成的 Hamilton 向量场从 $L$ 出发、终止于 $L'$ 的 time-$1$ 轨道：
\begin{equation}\label{eq:floergenerators}
\cX(L,L')=\{\gamma:[0,1]\to M\,|\,\gamma(0)\in L,\ \gamma(1)\in L',\ \dot\gamma(t)=X_{\varepsilon h}(\gamma(t))\},
\end{equation}
等价地，是在流作用下相互匹配的 $L$ 与 $L'$ 中的点对
\[
\cX(L,L')\simeq\{(p,p')\in L\times L'\,|\,\phi^1_{\varepsilon h}(p)=p'\},
\]
或者更简单地，是 $L^+\cap L'$ 中的点。（我们滥用记号，把 $\cX(L,L')$ 的元素不加区分地看作点、点对或 $X_{\varepsilon h}$ 的轨道。）注意，当 $L,L'$ 都是向量丛型时，$\cX(L,L')$ 在 $M$ 的每个节点处总包含一个生成元。

态射空间定义为
\begin{equation}\label{eq:morphisms}
\hom_{\F(M)}((L,\cE),(L',\cE'))=CF^*((L,\cE),(L',\cE');\varepsilon h)=\bigoplus_{(p,p')\in\cX(L,L')}\cE_p^*\otimes\cE'_{p'}.
\end{equation}
（另一种做法是把 $\F(M)$ 定义为这样一个有向范畴：对象是 $(L,\cE)$ 在正 Hamilton 流下的像，并对延拓元素 $e_{(L,\cE),\varepsilon}\in CF^*(\phi^1_{\varepsilon h}(L,\cE),(L,\cE))$ 作局部化；这与近期某些文献 \cite{AA,AS} 更为一致，但在我们的设定中这样做并无好处。）

在节点之外选取切丛 $TM$ 的平凡化，就确定了 $\F(M)$ 上的 $\Z$-分次。我们这里的优先选择是由各分量开层 $\C^*\subset M_e$ 上的径向线场所确定的平凡化。对象的分次则由 $TL$ 与所选线场之间角度的实值提升给出。这里同样有优先选择：对 v.b. 型对象，我们规定在每个分量中 $TL$ 与向外的径向线场在两端（$0$ 与 $\infty$）附近的角度为零；对点型对象，当 $L$ 是 $M_e$ 中以原点为圆心的圆时，规定该角度为 $-\pi/2$。按此约定，一切 Floer 上同调群都集中在 $0$ 次与 $1$ 次；对两个 v.b. 型对象而言，位于 $M$ 节点处的生成元在 $0$ 次。

\begin{remark}\label{rmk:notCY}
由于定义态射空间时用了正 Hamilton 扰动，当 $M$ 紧（所有分量都是 $\PP^1$）时，范畴 $\F(M)$ 从不是 Calabi--Yau 的。$\F(M)$ 的开--闭与闭--开映射的研究超出本文范围，但我们指出：$\F(M)$ 的 Hochschild 上同调应当同构于小正 Hamilton 函数的不动点 Floer 上同调，其同构由闭--开映射
\[
\mathcal{CO}:HF^*(\phi^1_{\varepsilon h})\to HH^*(\F(M))
\]
给出。例如，当 $M$ 由 $3g-3$ 条 $\PP^1$ 在 $2g-2$ 个三重点处相交组成时，存在正 Hamilton 函数，它有 $2g-2$ 个极小点（在节点处）、$3g-3$ 个鞍点与 $3g-3$ 个极大点。$CF^*(\phi^1_{\varepsilon h})$ 上的 Floer 微分与 $M$ 各分量内的 Morse 微分一致，于是每个极小点映为三个鞍点之和，并且
\[
\dim HF^0(\phi^1_{\varepsilon h})=1,\qquad
\dim HF^1(\phi^1_{\varepsilon h})=g,\qquad
\dim HF^2(\phi^1_{\varepsilon h})=3g-3,
\]
这与亏格 $g$ 曲线导出范畴的 Hochschild 上同调一致。
\end{remark}
% [Tutor Note: 易错] 此处「从不是 Calabi-Yau」指的是：正 Hamilton 扰动使对象在无穷远处被推挤，Floer 上同调的乘积不再有 Calabi-Yau 情形下的 Poincaré 对偶/单位元性质；这与「F(M) 的 Hochschild 上同调等于闭弦 Floer 上同调」并不矛盾，后者是开--闭对应，而非 CY 性质。依据：原文 Remark 3.4（本工程 03_术语表）。
